% Options for packages loaded elsewhere
\PassOptionsToPackage{unicode}{hyperref}
\PassOptionsToPackage{hyphens}{url}
%
\documentclass[
  brazilian,
]{article}
\usepackage{lmodern}
\usepackage{amssymb,amsmath}
\usepackage{ifxetex,ifluatex}
\ifnum 0\ifxetex 1\fi\ifluatex 1\fi=0 % if pdftex
  \usepackage[T1]{fontenc}
  \usepackage[utf8]{inputenc}
  \usepackage{textcomp} % provide euro and other symbols
\else % if luatex or xetex
  \usepackage{unicode-math}
  \defaultfontfeatures{Scale=MatchLowercase}
  \defaultfontfeatures[\rmfamily]{Ligatures=TeX,Scale=1}
\fi
% Use upquote if available, for straight quotes in verbatim environments
\IfFileExists{upquote.sty}{\usepackage{upquote}}{}
\IfFileExists{microtype.sty}{% use microtype if available
  \usepackage[]{microtype}
  \UseMicrotypeSet[protrusion]{basicmath} % disable protrusion for tt fonts
}{}
\makeatletter
\@ifundefined{KOMAClassName}{% if non-KOMA class
  \IfFileExists{parskip.sty}{%
    \usepackage{parskip}
  }{% else
    \setlength{\parindent}{0pt}
    \setlength{\parskip}{6pt plus 2pt minus 1pt}}
}{% if KOMA class
  \KOMAoptions{parskip=half}}
\makeatother
\usepackage{xcolor}
\IfFileExists{xurl.sty}{\usepackage{xurl}}{} % add URL line breaks if available
\IfFileExists{bookmark.sty}{\usepackage{bookmark}}{\usepackage{hyperref}}
\hypersetup{
  pdftitle={Parte III: Processamento de Linguagem Natural para Linguistas},
  pdfauthor={CiberExt 26-29 · FEELT38103 · Universidade Federal de Uberlândia},
  pdflang={pt-BR},
  hidelinks,
  pdfcreator={LaTeX via pandoc}}
\urlstyle{same} % disable monospaced font for URLs
\usepackage{color}
\usepackage{fancyvrb}
\newcommand{\VerbBar}{|}
\newcommand{\VERB}{\Verb[commandchars=\\\{\}]}
\DefineVerbatimEnvironment{Highlighting}{Verbatim}{commandchars=\\\{\}}
% Add ',fontsize=\small' for more characters per line
\newenvironment{Shaded}{}{}
\newcommand{\AlertTok}[1]{\textcolor[rgb]{1.00,0.00,0.00}{\textbf{#1}}}
\newcommand{\AnnotationTok}[1]{\textcolor[rgb]{0.38,0.63,0.69}{\textbf{\textit{#1}}}}
\newcommand{\AttributeTok}[1]{\textcolor[rgb]{0.49,0.56,0.16}{#1}}
\newcommand{\BaseNTok}[1]{\textcolor[rgb]{0.25,0.63,0.44}{#1}}
\newcommand{\BuiltInTok}[1]{#1}
\newcommand{\CharTok}[1]{\textcolor[rgb]{0.25,0.44,0.63}{#1}}
\newcommand{\CommentTok}[1]{\textcolor[rgb]{0.38,0.63,0.69}{\textit{#1}}}
\newcommand{\CommentVarTok}[1]{\textcolor[rgb]{0.38,0.63,0.69}{\textbf{\textit{#1}}}}
\newcommand{\ConstantTok}[1]{\textcolor[rgb]{0.53,0.00,0.00}{#1}}
\newcommand{\ControlFlowTok}[1]{\textcolor[rgb]{0.00,0.44,0.13}{\textbf{#1}}}
\newcommand{\DataTypeTok}[1]{\textcolor[rgb]{0.56,0.13,0.00}{#1}}
\newcommand{\DecValTok}[1]{\textcolor[rgb]{0.25,0.63,0.44}{#1}}
\newcommand{\DocumentationTok}[1]{\textcolor[rgb]{0.73,0.13,0.13}{\textit{#1}}}
\newcommand{\ErrorTok}[1]{\textcolor[rgb]{1.00,0.00,0.00}{\textbf{#1}}}
\newcommand{\ExtensionTok}[1]{#1}
\newcommand{\FloatTok}[1]{\textcolor[rgb]{0.25,0.63,0.44}{#1}}
\newcommand{\FunctionTok}[1]{\textcolor[rgb]{0.02,0.16,0.49}{#1}}
\newcommand{\ImportTok}[1]{#1}
\newcommand{\InformationTok}[1]{\textcolor[rgb]{0.38,0.63,0.69}{\textbf{\textit{#1}}}}
\newcommand{\KeywordTok}[1]{\textcolor[rgb]{0.00,0.44,0.13}{\textbf{#1}}}
\newcommand{\NormalTok}[1]{#1}
\newcommand{\OperatorTok}[1]{\textcolor[rgb]{0.40,0.40,0.40}{#1}}
\newcommand{\OtherTok}[1]{\textcolor[rgb]{0.00,0.44,0.13}{#1}}
\newcommand{\PreprocessorTok}[1]{\textcolor[rgb]{0.74,0.48,0.00}{#1}}
\newcommand{\RegionMarkerTok}[1]{#1}
\newcommand{\SpecialCharTok}[1]{\textcolor[rgb]{0.25,0.44,0.63}{#1}}
\newcommand{\SpecialStringTok}[1]{\textcolor[rgb]{0.73,0.40,0.53}{#1}}
\newcommand{\StringTok}[1]{\textcolor[rgb]{0.25,0.44,0.63}{#1}}
\newcommand{\VariableTok}[1]{\textcolor[rgb]{0.10,0.09,0.49}{#1}}
\newcommand{\VerbatimStringTok}[1]{\textcolor[rgb]{0.25,0.44,0.63}{#1}}
\newcommand{\WarningTok}[1]{\textcolor[rgb]{0.38,0.63,0.69}{\textbf{\textit{#1}}}}
\setlength{\emergencystretch}{3em} % prevent overfull lines
\providecommand{\tightlist}{%
  \setlength{\itemsep}{0pt}\setlength{\parskip}{0pt}}
\setcounter{secnumdepth}{-\maxdimen} % remove section numbering
\ifxetex
  % Load polyglossia as late as possible: uses bidi with RTL langages (e.g. Hebrew, Arabic)
  \usepackage{polyglossia}
  \setmainlanguage[]{brazilian}
\else
  \usepackage[shorthands=off,main=brazilian]{babel}
\fi

\title{Parte III: Processamento de Linguagem Natural para Linguistas}
\usepackage{etoolbox}
\makeatletter
\providecommand{\subtitle}[1]{% add subtitle to \maketitle
  \apptocmd{\@title}{\par {\large #1 \par}}{}{}
}
\makeatother
\subtitle{Unidade 3 - Respostas dos Exercícios}
\author{CiberExt 26-29 · FEELT38103 · Universidade Federal de
Uberlândia}
\date{Agosto de 2026}

\begin{document}
\maketitle

\hypertarget{respostas-encontrando-padruxf5es-linguuxedsticos-com-o-spacy}{%
\section{Respostas --- Encontrando padrões linguísticos com o
spaCy}\label{respostas-encontrando-padruxf5es-linguuxedsticos-com-o-spacy}}

\hypertarget{parte-a-conceitos}{%
\subsection{Parte A --- Conceitos}\label{parte-a-conceitos}}

\textbf{Resposta 1.}

\begin{itemize}
\tightlist
\item
  O \emph{Matcher} casa \textbf{sequências lineares de \emph{Tokens}}
  com base em seus atributos (classe gramatical, traços morfológicos,
  forma, lema etc.). É adequado, por exemplo, para localizar sequências
  ``pronome + verbo''.
\item
  O \emph{DependencyMatcher} casa \textbf{padrões na árvore de
  dependências}, isto é, relações de cabeça e dependente. É adequado
  para localizar, por exemplo, todos os sujeitos nominais de verbos,
  independentemente da ordem linear das palavras.
\item
  O \emph{PhraseMatcher} casa \textbf{objetos \emph{Doc} contra objetos
  \emph{Doc}} e é otimizado para listas grandes de termos fixos. É
  adequado, por exemplo, para localizar milhares de nomes de
  medicamentos ou de topônimos em um corpus.
\end{itemize}

\textbf{Resposta 2.}

A chave \texttt{OP} define a \textbf{quantificação} de um padrão de
\emph{Token}:

\begin{itemize}
\tightlist
\item
  \texttt{!} --- nega o padrão; ele deve ocorrer exatamente zero vezes.
\item
  \texttt{?} --- torna o padrão opcional; pode ocorrer zero ou uma vez.
\item
  \texttt{+} --- exige uma ou mais ocorrências.
\item
  \texttt{*} --- permite zero ou mais ocorrências.
\end{itemize}

O operador \texttt{?} é preferível ao \texttt{+} quando o elemento é
\textbf{facultativo, mas não repetível}, como o determinante em um
sintagma nominal do inglês:
\texttt{{[}\{\textquotesingle{}POS\textquotesingle{}:\ \textquotesingle{}DET\textquotesingle{},\ \textquotesingle{}OP\textquotesingle{}:\ \textquotesingle{}?\textquotesingle{}\},\ \{\textquotesingle{}POS\textquotesingle{}:\ \textquotesingle{}NOUN\textquotesingle{}\}{]}}
casa tanto \emph{the capital} quanto \emph{capitals}, mas não permitiria
uma sequência de determinantes.

\textbf{Resposta 3.}

\texttt{POS} refere-se à etiqueta \textbf{geral} (\emph{coarse-grained})
de classe gramatical, definida pelo conjunto universal das Dependências
Universais (\texttt{NOUN}, \texttt{VERB}, \texttt{ADJ}, \texttt{PROPN}
etc.). \texttt{TAG} refere-se à etiqueta \textbf{refinada}
(\emph{fine-grained}), específica do corpus de treinamento do modelo ---
no caso do inglês, o conjunto do Penn Treebank (\texttt{NN},
\texttt{NNS}, \texttt{NNP}, \texttt{VBD} etc.).

\texttt{NNP} é uma etiqueta do Penn Treebank e, portanto, só é válida
como valor de \texttt{TAG}. O equivalente geral seria \texttt{PROPN} em
\texttt{POS}, mas \texttt{NNP} é mais específico, pois distingue nome
próprio singular de plural (\texttt{NNPS}).

\textbf{Resposta 4.}

Quando o valor de \texttt{MORPH} é uma \emph{string}, o spaCy exige uma
correspondência \textbf{exata} de todo o conjunto de traços
morfológicos. Com \texttt{IS\_SUPERSET}, o spaCy verifica se o conjunto
de traços do \emph{Token} \textbf{contém} o conjunto fornecido no
padrão, ou seja, se é um superconjunto dele.

Isso é necessário porque quase nunca conhecemos de antemão a combinação
completa de traços de um \emph{Token}. Ao buscar apenas verbos no
passado, por exemplo, queremos casar tanto
\texttt{Tense=Past\textbar{}VerbForm=Fin} quanto
\texttt{Aspect=Perf\textbar{}Tense=Past\textbar{}VerbForm=Part}; um
único traço, \texttt{Tense=Past}, basta como critério.

\textbf{Resposta 5.}

\texttt{RIGHT\_ID} é o \textbf{nome do padrão atual}, usado para que
padrões posteriores possam referenciá-lo. \texttt{LEFT\_ID} é a
\textbf{referência a um padrão já definido}, ao qual o padrão atual se
conecta por meio do operador relacional em \texttt{REL\_OP}.

O primeiro dicionário não possui \texttt{LEFT\_ID} porque é a
\textbf{âncora}: não há nenhum padrão anterior ao qual ele possa se
ligar. A âncora é o ponto de partida da corrente.

\textbf{Resposta 6.}

Um objeto \emph{Span} representa uma \textbf{fatia contígua} de um
objeto \emph{Doc}, definida por um índice inicial e um final.
Correspondências de dependência, porém, ligam \emph{Tokens} que podem
estar arbitrariamente distantes na ordem linear e com material
interveniente que não faz parte da correspondência (por exemplo, em
\emph{the movement in New York began}, o sujeito \texttt{movement} e o
verbo \texttt{began} não são adjacentes). Como não formam uma fatia
contígua, não podem ser representados por um \emph{Span}; por isso o
spaCy devolve os índices individuais dos \emph{Tokens}.

\hypertarget{parte-b-programauxe7uxe3o}{%
\subsection{Parte B --- Programação}\label{parte-b-programauxe7uxe3o}}

\textbf{Resposta 7.}

\begin{Shaded}
\begin{Highlighting}[]
\ImportTok{from}\NormalTok{ spacy.matcher }\ImportTok{import}\NormalTok{ Matcher}

\CommentTok{\# Cria um novo Matcher com o vocabulário do modelo}
\NormalTok{adj\_matcher }\OperatorTok{=}\NormalTok{ Matcher(vocab}\OperatorTok{=}\NormalTok{nlp.vocab)}

\CommentTok{\# Define o padrão: um adjetivo seguido de um substantivo}
\NormalTok{adj\_noun }\OperatorTok{=}\NormalTok{ [\{}\StringTok{\textquotesingle{}POS\textquotesingle{}}\NormalTok{: }\StringTok{\textquotesingle{}ADJ\textquotesingle{}}\NormalTok{\}, \{}\StringTok{\textquotesingle{}POS\textquotesingle{}}\NormalTok{: }\StringTok{\textquotesingle{}NOUN\textquotesingle{}}\NormalTok{\}]}

\CommentTok{\# Adiciona o padrão ao Matcher}
\NormalTok{adj\_matcher.add(}\StringTok{\textquotesingle{}adj+noun\textquotesingle{}}\NormalTok{, patterns}\OperatorTok{=}\NormalTok{[adj\_noun])}

\CommentTok{\# Aplica o Matcher ao Doc e obtém os resultados como Spans}
\NormalTok{adj\_results }\OperatorTok{=}\NormalTok{ adj\_matcher(doc, as\_spans}\OperatorTok{=}\VariableTok{True}\NormalTok{)}

\CommentTok{\# Imprime as 20 primeiras correspondências}
\ControlFlowTok{for}\NormalTok{ result }\KeywordTok{in}\NormalTok{ adj\_results[:}\DecValTok{20}\NormalTok{]:}
    \BuiltInTok{print}\NormalTok{(result)}
\end{Highlighting}
\end{Shaded}

\textbf{Resposta 8.}

\begin{Shaded}
\begin{Highlighting}[]
\CommentTok{\# Permite um ou mais adjetivos antes do substantivo}
\NormalTok{adj\_noun\_plus }\OperatorTok{=}\NormalTok{ [\{}\StringTok{\textquotesingle{}POS\textquotesingle{}}\NormalTok{: }\StringTok{\textquotesingle{}ADJ\textquotesingle{}}\NormalTok{, }\StringTok{\textquotesingle{}OP\textquotesingle{}}\NormalTok{: }\StringTok{\textquotesingle{}+\textquotesingle{}}\NormalTok{\}, \{}\StringTok{\textquotesingle{}POS\textquotesingle{}}\NormalTok{: }\StringTok{\textquotesingle{}NOUN\textquotesingle{}}\NormalTok{\}]}

\CommentTok{\# Versão sem \textquotesingle{}greedy\textquotesingle{}}
\NormalTok{matcher\_a }\OperatorTok{=}\NormalTok{ Matcher(vocab}\OperatorTok{=}\NormalTok{nlp.vocab)}
\NormalTok{matcher\_a.add(}\StringTok{\textquotesingle{}adj+noun\textquotesingle{}}\NormalTok{, patterns}\OperatorTok{=}\NormalTok{[adj\_noun\_plus])}
\NormalTok{results\_a }\OperatorTok{=}\NormalTok{ matcher\_a(doc, as\_spans}\OperatorTok{=}\VariableTok{True}\NormalTok{)}

\CommentTok{\# Versão com \textquotesingle{}greedy\textquotesingle{}}
\NormalTok{matcher\_b }\OperatorTok{=}\NormalTok{ Matcher(vocab}\OperatorTok{=}\NormalTok{nlp.vocab)}
\NormalTok{matcher\_b.add(}\StringTok{\textquotesingle{}adj+noun\textquotesingle{}}\NormalTok{, patterns}\OperatorTok{=}\NormalTok{[adj\_noun\_plus], greedy}\OperatorTok{=}\StringTok{\textquotesingle{}LONGEST\textquotesingle{}}\NormalTok{)}
\NormalTok{results\_b }\OperatorTok{=}\NormalTok{ matcher\_b(doc, as\_spans}\OperatorTok{=}\VariableTok{True}\NormalTok{)}

\BuiltInTok{print}\NormalTok{(}\StringTok{\textquotesingle{}Sem greedy:\textquotesingle{}}\NormalTok{, }\BuiltInTok{len}\NormalTok{(results\_a))}
\BuiltInTok{print}\NormalTok{(}\StringTok{\textquotesingle{}Com greedy:\textquotesingle{}}\NormalTok{, }\BuiltInTok{len}\NormalTok{(results\_b))}
\end{Highlighting}
\end{Shaded}

\textbf{Explicação.} Sem \texttt{greedy}, o operador \texttt{+} devolve
uma correspondência a \emph{cada} posição em que o padrão pode ser
satisfeito: para a sequência \emph{large public space}, o spaCy devolve
tanto \texttt{public\ space} quanto \texttt{large\ public\ space}. Com
\texttt{greedy=\textquotesingle{}LONGEST\textquotesingle{}}, o spaCy
filtra as correspondências sobrepostas e mantém apenas a mais longa de
cada grupo, reduzindo o total de resultados.

\textbf{Resposta 9.}

\begin{Shaded}
\begin{Highlighting}[]
\CommentTok{\# Cria um Matcher para traços morfológicos}
\NormalTok{plural\_matcher }\OperatorTok{=}\NormalTok{ Matcher(vocab}\OperatorTok{=}\NormalTok{nlp.vocab)}

\CommentTok{\# Define o padrão: substantivos cujos traços contenham Number=Plur}
\NormalTok{plural\_noun }\OperatorTok{=}\NormalTok{ [\{}\StringTok{\textquotesingle{}POS\textquotesingle{}}\NormalTok{: }\StringTok{\textquotesingle{}NOUN\textquotesingle{}}\NormalTok{, }\StringTok{\textquotesingle{}MORPH\textquotesingle{}}\NormalTok{: \{}\StringTok{\textquotesingle{}IS\_SUPERSET\textquotesingle{}}\NormalTok{: [}\StringTok{\textquotesingle{}Number=Plur\textquotesingle{}}\NormalTok{]\}\}]}

\CommentTok{\# Adiciona o padrão ao Matcher}
\NormalTok{plural\_matcher.add(}\StringTok{\textquotesingle{}plural\_noun\textquotesingle{}}\NormalTok{, patterns}\OperatorTok{=}\NormalTok{[plural\_noun])}

\CommentTok{\# Aplica o Matcher ao Doc}
\NormalTok{plural\_results }\OperatorTok{=}\NormalTok{ plural\_matcher(doc, as\_spans}\OperatorTok{=}\VariableTok{True}\NormalTok{)}

\CommentTok{\# Imprime as 15 primeiras correspondências e seus traços morfológicos}
\ControlFlowTok{for}\NormalTok{ result }\KeywordTok{in}\NormalTok{ plural\_results[:}\DecValTok{15}\NormalTok{]:}
    \BuiltInTok{print}\NormalTok{(result, }\StringTok{\textquotesingle{}}\CharTok{\textbackslash{}t}\StringTok{\textquotesingle{}}\NormalTok{, result[}\DecValTok{0}\NormalTok{].morph)}
\end{Highlighting}
\end{Shaded}

\textbf{Resposta 10.}

\begin{Shaded}
\begin{Highlighting}[]
\ImportTok{from}\NormalTok{ spacy.matcher }\ImportTok{import}\NormalTok{ DependencyMatcher}

\CommentTok{\# Cria um novo DependencyMatcher}
\NormalTok{obj\_matcher }\OperatorTok{=}\NormalTok{ DependencyMatcher(vocab}\OperatorTok{=}\NormalTok{nlp.vocab)}

\CommentTok{\# Define a corrente: âncora (verbo) e seu objeto direto}
\NormalTok{verb\_dobj }\OperatorTok{=}\NormalTok{ [\{}\StringTok{\textquotesingle{}RIGHT\_ID\textquotesingle{}}\NormalTok{: }\StringTok{\textquotesingle{}verb\textquotesingle{}}\NormalTok{, }\StringTok{\textquotesingle{}RIGHT\_ATTRS\textquotesingle{}}\NormalTok{: \{}\StringTok{\textquotesingle{}POS\textquotesingle{}}\NormalTok{: }\StringTok{\textquotesingle{}VERB\textquotesingle{}}\NormalTok{\}\},}
\NormalTok{             \{}\StringTok{\textquotesingle{}LEFT\_ID\textquotesingle{}}\NormalTok{: }\StringTok{\textquotesingle{}verb\textquotesingle{}}\NormalTok{, }\StringTok{\textquotesingle{}REL\_OP\textquotesingle{}}\NormalTok{: }\StringTok{\textquotesingle{}\textgreater{}\textquotesingle{}}\NormalTok{, }\StringTok{\textquotesingle{}RIGHT\_ID\textquotesingle{}}\NormalTok{: }\StringTok{\textquotesingle{}d\_object\textquotesingle{}}\NormalTok{,}
              \StringTok{\textquotesingle{}RIGHT\_ATTRS\textquotesingle{}}\NormalTok{: \{}\StringTok{\textquotesingle{}DEP\textquotesingle{}}\NormalTok{: }\StringTok{\textquotesingle{}dobj\textquotesingle{}}\NormalTok{\}\}}
\NormalTok{            ]}

\CommentTok{\# Adiciona o padrão ao matcher}
\NormalTok{obj\_matcher.add(}\StringTok{\textquotesingle{}verb\_dobj\textquotesingle{}}\NormalTok{, patterns}\OperatorTok{=}\NormalTok{[verb\_dobj])}

\CommentTok{\# Aplica o matcher ao Doc}
\NormalTok{obj\_matches }\OperatorTok{=}\NormalTok{ obj\_matcher(doc)}

\CommentTok{\# Imprime as 20 primeiras correspondências}
\ControlFlowTok{for}\NormalTok{ match }\KeywordTok{in}\NormalTok{ obj\_matches[:}\DecValTok{20}\NormalTok{]:}

    \CommentTok{\# Desempacota a tupla em nome do padrão e lista de índices}
\NormalTok{    pattern\_name, matches }\OperatorTok{=}\NormalTok{ match[}\DecValTok{0}\NormalTok{], match[}\DecValTok{1}\NormalTok{]}

    \CommentTok{\# Os índices seguem a ordem em que os padrões foram definidos}
\NormalTok{    verb, dobject }\OperatorTok{=}\NormalTok{ matches[}\DecValTok{0}\NormalTok{], matches[}\DecValTok{1}\NormalTok{]}

    \BuiltInTok{print}\NormalTok{(nlp.vocab[pattern\_name].text, }\StringTok{\textquotesingle{}}\CharTok{\textbackslash{}t}\StringTok{\textquotesingle{}}\NormalTok{, doc[verb], }\StringTok{\textquotesingle{}...\textquotesingle{}}\NormalTok{, doc[dobject])}
\end{Highlighting}
\end{Shaded}

\textbf{Resposta 11.}

\begin{Shaded}
\begin{Highlighting}[]
\CommentTok{\# O terceiro elo parte de \textquotesingle{}d\_object\textquotesingle{}, e não da âncora \textquotesingle{}verb\textquotesingle{}}
\NormalTok{verb\_dobj\_det }\OperatorTok{=}\NormalTok{ [\{}\StringTok{\textquotesingle{}RIGHT\_ID\textquotesingle{}}\NormalTok{: }\StringTok{\textquotesingle{}verb\textquotesingle{}}\NormalTok{, }\StringTok{\textquotesingle{}RIGHT\_ATTRS\textquotesingle{}}\NormalTok{: \{}\StringTok{\textquotesingle{}POS\textquotesingle{}}\NormalTok{: }\StringTok{\textquotesingle{}VERB\textquotesingle{}}\NormalTok{\}\},}
\NormalTok{                 \{}\StringTok{\textquotesingle{}LEFT\_ID\textquotesingle{}}\NormalTok{: }\StringTok{\textquotesingle{}verb\textquotesingle{}}\NormalTok{, }\StringTok{\textquotesingle{}REL\_OP\textquotesingle{}}\NormalTok{: }\StringTok{\textquotesingle{}\textgreater{}\textquotesingle{}}\NormalTok{, }\StringTok{\textquotesingle{}RIGHT\_ID\textquotesingle{}}\NormalTok{: }\StringTok{\textquotesingle{}d\_object\textquotesingle{}}\NormalTok{,}
                  \StringTok{\textquotesingle{}RIGHT\_ATTRS\textquotesingle{}}\NormalTok{: \{}\StringTok{\textquotesingle{}DEP\textquotesingle{}}\NormalTok{: }\StringTok{\textquotesingle{}dobj\textquotesingle{}}\NormalTok{\}\},}
\NormalTok{                 \{}\StringTok{\textquotesingle{}LEFT\_ID\textquotesingle{}}\NormalTok{: }\StringTok{\textquotesingle{}d\_object\textquotesingle{}}\NormalTok{, }\StringTok{\textquotesingle{}REL\_OP\textquotesingle{}}\NormalTok{: }\StringTok{\textquotesingle{}\textgreater{}\textquotesingle{}}\NormalTok{, }\StringTok{\textquotesingle{}RIGHT\_ID\textquotesingle{}}\NormalTok{: }\StringTok{\textquotesingle{}determiner\textquotesingle{}}\NormalTok{,}
                  \StringTok{\textquotesingle{}RIGHT\_ATTRS\textquotesingle{}}\NormalTok{: \{}\StringTok{\textquotesingle{}DEP\textquotesingle{}}\NormalTok{: }\StringTok{\textquotesingle{}det\textquotesingle{}}\NormalTok{\}\}}
\NormalTok{                ]}

\CommentTok{\# Cria um matcher novo para evitar misturar com padrões anteriores}
\NormalTok{det\_matcher }\OperatorTok{=}\NormalTok{ DependencyMatcher(vocab}\OperatorTok{=}\NormalTok{nlp.vocab)}
\NormalTok{det\_matcher.add(}\StringTok{\textquotesingle{}verb\_dobj\_det\textquotesingle{}}\NormalTok{, patterns}\OperatorTok{=}\NormalTok{[verb\_dobj\_det])}

\NormalTok{det\_matches }\OperatorTok{=}\NormalTok{ det\_matcher(doc)}

\ControlFlowTok{for}\NormalTok{ match }\KeywordTok{in}\NormalTok{ det\_matches[:}\DecValTok{20}\NormalTok{]:}

\NormalTok{    pattern\_name, matches }\OperatorTok{=}\NormalTok{ match[}\DecValTok{0}\NormalTok{], match[}\DecValTok{1}\NormalTok{]}

    \CommentTok{\# A ordem dos índices segue a ordem de definição dos padrões}
\NormalTok{    verb, dobject, determiner }\OperatorTok{=}\NormalTok{ matches[}\DecValTok{0}\NormalTok{], matches[}\DecValTok{1}\NormalTok{], matches[}\DecValTok{2}\NormalTok{]}

    \BuiltInTok{print}\NormalTok{(nlp.vocab[pattern\_name].text, }\StringTok{\textquotesingle{}}\CharTok{\textbackslash{}t}\StringTok{\textquotesingle{}}\NormalTok{,}
\NormalTok{          doc[verb], }\StringTok{\textquotesingle{}...\textquotesingle{}}\NormalTok{, doc[determiner], doc[dobject])}
\end{Highlighting}
\end{Shaded}

\textbf{Explicação.} O determinante é dependente do
\textbf{substantivo}, não do verbo. Por isso \texttt{LEFT\_ID} deve ser
\texttt{d\_object}: estamos prolongando a corrente para a direita a
partir do objeto direto, e não abrindo uma nova corrente a partir da
âncora.

\textbf{Resposta 12.}

\begin{Shaded}
\begin{Highlighting}[]
\ImportTok{from}\NormalTok{ wasabi }\ImportTok{import}\NormalTok{ Printer}

\CommentTok{\# Inicializa o Printer}
\NormalTok{match }\OperatorTok{=}\NormalTok{ Printer()}

\CommentTok{\# Percorre as correspondências mantendo a contagem}
\ControlFlowTok{for}\NormalTok{ i, result }\KeywordTok{in} \BuiltInTok{enumerate}\NormalTok{(adj\_results, start}\OperatorTok{=}\DecValTok{1}\NormalTok{):}

    \BuiltInTok{print}\NormalTok{(i,}
\NormalTok{          doc[result.start }\OperatorTok{{-}} \DecValTok{5}\NormalTok{: result.start],   }\CommentTok{\# contexto à esquerda}
\NormalTok{          match.text(result, color}\OperatorTok{=}\StringTok{\textquotesingle{}blue\textquotesingle{}}\NormalTok{, no\_print}\OperatorTok{=}\VariableTok{True}\NormalTok{),  }\CommentTok{\# correspondência em azul}
\NormalTok{          doc[result.end: result.end }\OperatorTok{+} \DecValTok{5}\NormalTok{]        }\CommentTok{\# contexto à direita}
\NormalTok{         )}
\end{Highlighting}
\end{Shaded}

\hypertarget{parte-c-reflexuxe3o}{%
\subsection{Parte C --- Reflexão}\label{parte-c-reflexuxe3o}}

\textbf{Resposta 13.}

O texto de origem contém quebras de linha (\texttt{\textbackslash{}n})
que o spaCy preserva como \emph{Tokens}. Quando uma dessas quebras cai
na janela de contexto, ela é impressa literalmente e desloca a saída
para outra linha.

Uma solução é substituir os caracteres de espaçamento pelo caractere de
espaço antes de imprimir:

\begin{Shaded}
\begin{Highlighting}[]
\ControlFlowTok{for}\NormalTok{ i, result }\KeywordTok{in} \BuiltInTok{enumerate}\NormalTok{(adj\_results, start}\OperatorTok{=}\DecValTok{1}\NormalTok{):}

    \CommentTok{\# Converte as fatias em texto e normaliza os espaços em branco}
\NormalTok{    left }\OperatorTok{=} \StringTok{\textquotesingle{} \textquotesingle{}}\NormalTok{.join(doc[result.start }\OperatorTok{{-}} \DecValTok{5}\NormalTok{: result.start].text.split())}
\NormalTok{    right }\OperatorTok{=} \StringTok{\textquotesingle{} \textquotesingle{}}\NormalTok{.join(doc[result.end: result.end }\OperatorTok{+} \DecValTok{5}\NormalTok{].text.split())}

    \BuiltInTok{print}\NormalTok{(i, left, match.text(result, color}\OperatorTok{=}\StringTok{\textquotesingle{}blue\textquotesingle{}}\NormalTok{, no\_print}\OperatorTok{=}\VariableTok{True}\NormalTok{), right)}
\end{Highlighting}
\end{Shaded}

O método \texttt{split()} sem argumentos divide a \emph{string} em
qualquer sequência de espaços em branco --- incluindo
\texttt{\textbackslash{}n} e \texttt{\textbackslash{}t} --- e
\texttt{join()} a reconstitui com espaços simples.

Alternativamente, poderíamos limpar o texto antes de fornecê-lo ao
modelo de linguagem, ou filtrar os \emph{Tokens} cujo atributo
\texttt{is\_space} seja \texttt{True}.

\textbf{Resposta 14.}

A construção envolve uma relação \textbf{sintática}, e não meramente
linear: o sintagma nominal é o objeto da preposição, que por sua vez é
dependente do verbo. Palavras podem intervir entre os elementos
(\emph{depend heavily on the government}), o que torna o \emph{Matcher}
linear frágil. Portanto, a escolha adequada é o
\textbf{DependencyMatcher}.

A regra teria três elos:

\begin{Shaded}
\begin{Highlighting}[]
\NormalTok{verb\_prep\_np }\OperatorTok{=}\NormalTok{ [}
    \CommentTok{\# Âncora: o verbo}
\NormalTok{    \{}\StringTok{\textquotesingle{}RIGHT\_ID\textquotesingle{}}\NormalTok{: }\StringTok{\textquotesingle{}verb\textquotesingle{}}\NormalTok{, }\StringTok{\textquotesingle{}RIGHT\_ATTRS\textquotesingle{}}\NormalTok{: \{}\StringTok{\textquotesingle{}POS\textquotesingle{}}\NormalTok{: }\StringTok{\textquotesingle{}VERB\textquotesingle{}}\NormalTok{\}\},}

    \CommentTok{\# A preposição é dependente do verbo, com a relação \textquotesingle{}prep\textquotesingle{}}
\NormalTok{    \{}\StringTok{\textquotesingle{}LEFT\_ID\textquotesingle{}}\NormalTok{: }\StringTok{\textquotesingle{}verb\textquotesingle{}}\NormalTok{, }\StringTok{\textquotesingle{}REL\_OP\textquotesingle{}}\NormalTok{: }\StringTok{\textquotesingle{}\textgreater{}\textquotesingle{}}\NormalTok{, }\StringTok{\textquotesingle{}RIGHT\_ID\textquotesingle{}}\NormalTok{: }\StringTok{\textquotesingle{}preposition\textquotesingle{}}\NormalTok{,}
     \StringTok{\textquotesingle{}RIGHT\_ATTRS\textquotesingle{}}\NormalTok{: \{}\StringTok{\textquotesingle{}DEP\textquotesingle{}}\NormalTok{: }\StringTok{\textquotesingle{}prep\textquotesingle{}}\NormalTok{\}\},}

    \CommentTok{\# O objeto da preposição é dependente da preposição, com a relação \textquotesingle{}pobj\textquotesingle{}}
\NormalTok{    \{}\StringTok{\textquotesingle{}LEFT\_ID\textquotesingle{}}\NormalTok{: }\StringTok{\textquotesingle{}preposition\textquotesingle{}}\NormalTok{, }\StringTok{\textquotesingle{}REL\_OP\textquotesingle{}}\NormalTok{: }\StringTok{\textquotesingle{}\textgreater{}\textquotesingle{}}\NormalTok{, }\StringTok{\textquotesingle{}RIGHT\_ID\textquotesingle{}}\NormalTok{: }\StringTok{\textquotesingle{}prep\_object\textquotesingle{}}\NormalTok{,}
     \StringTok{\textquotesingle{}RIGHT\_ATTRS\textquotesingle{}}\NormalTok{: \{}\StringTok{\textquotesingle{}DEP\textquotesingle{}}\NormalTok{: }\StringTok{\textquotesingle{}pobj\textquotesingle{}}\NormalTok{\}\}}
\NormalTok{]}
\end{Highlighting}
\end{Shaded}

\textbf{Justificativa das escolhas:}

\begin{itemize}
\tightlist
\item
  A âncora é o verbo, porque é o elemento em torno do qual a construção
  se organiza.
\item
  O segundo elo parte do verbo
  (\texttt{LEFT\_ID:\ \textquotesingle{}verb\textquotesingle{}}), pois a
  preposição é dependente dele.
\item
  O terceiro elo parte da preposição
  (\texttt{LEFT\_ID:\ \textquotesingle{}preposition\textquotesingle{}}),
  e \textbf{não} do verbo, porque o sintagma nominal é dependente da
  preposição --- trata-se de um prolongamento da mesma corrente.
\item
  Se quiséssemos restringir o verbo a um lema específico, bastaria
  acrescentar
  \texttt{\textquotesingle{}LEMMA\textquotesingle{}:\ \textquotesingle{}depend\textquotesingle{}}
  ao dicionário \texttt{RIGHT\_ATTRS} da âncora.
\end{itemize}

\end{document}
